% Complete corrected negation and arbitrary-order completeness proofs.

\section{The full Stavi completeness proof}
\label{sec:annex-p09-stavi-completeness}

\begin{annexitem}{annex:p9.stavi.completion.primitives}{thm:p09-stavi-completion-primitives}{Elementary simple formulas and all possible boundaries}
\begin{proof}
On actual \(a<b\), the closed interval \([a,b]\)
has exactly its two endpoint points precisely
when the open interval is empty. Two forced
true singletons therefore express \(E\).
A one-cell unflagged partition of \((a,b)\)
expresses nonemptiness and universal truth of
its label; adding \(E\) gives
\(\operatorname{All}_B\), including empty
intervals. A point label \(B(z)\) is the
basic one-point partition on \([z,z]\).
Thus \(\operatorname{Occ}_B\) is simple
by the proved positive existential closure.
Classical quantifier negation gives the
two displayed complement identities.

For \(j\geq1\), exact cardinality \(j\)
is the partition into \(j\) true-labelled
forced singletons; cardinality zero is
\(E\). At most \(j\) is their finite
disjunction. At least \(j+1\) is an
increasing tuple of \(j+1\) true-labelled
points, a positive existential formula
and hence simple. These give both
complements of each exact or bounded
cardinality condition. Order atoms
are handled by trichotomy, and a
negated unary label remains a
finite canonical temporal label.
None of this uses partition-negation
closure.

Let \(S\) be the nonempty occurrence
set of \(B\) in \((a,b)\). If its
infimum exists in \(T\), it lies
in \([a,b)\). If it is \(a\),
then at every later cut there is
an \(S\)-point strictly before
that cut; this is exactly
\(\partial^+B(a)\). Conversely
that formula supplies nonemptiness
by its instance at \(b\) and
gives infimum \(a\).

For \(a<c<b\), absence of \(B\)
on \((a,c)\) and \(B(c)\)
make \(c\) the minimum of \(S\).
If it is an unattained infimum,
then \(B(c)\) is false and
every later cut has an \(S\)-point
above \(c\) but below it.
This gives \(\partial^+B(c)\).
Conversely the absence and
this limit condition give a
nonempty set above \(c\)
with no larger lower bound,
so \(c\) is its infimum.
These prove the actual-point
descriptor, including an
attained point that also
has occurrences accumulating
immediately to its right.

If no infimum exists, let
\[
                        L=\{u:\forall s\in S,\ u<s\}.
\]
This set contains \(a\) and
its complement contains \(S\).
It is downward closed.
If it had a greatest point,
that point would be an
infimum, unless \(S\)
had a minimum; the latter
would itself be an infimum.
Both possibilities contradict
the assumption. If the
complement had a least
point \(c\), then \(c\leq s\)
for all \(s\in S\).
But \(c\notin L\), so
some \(s\leq c\), making
\(c\) a minimum of \(S\).
Thus this cut is proper.

No \(S\)-point lies to its
left. For a point \(r\)
on its right, some
\(s\in S\) satisfies
\(s\leq r\). If equality
holds, choose an earlier
member of \(S\), since
\(S\) has no minimum.
This proves that \(B\)
is coinitial immediately
after the gap. Hence it
is the \(\neg B\)-gap
succeeding \(a\), and
lies before some actual
occurrence below \(b\).
The gamma formula and
\(\operatorname{Occ}_B(a,b)\)
express exactly this case.
Conversely those conditions
give the same boundary
gap by the proved
uniqueness of the initial
run. This proves exhaustion
without assuming completeness.
Reversal gives the upper
classification and the
exact supremum descriptor
\[
 a<c<b\land\operatorname{All}_{\neg B}(c,b)
                              \land(B(c)\lor\partial^-B(c)).
\]
All limit tests here have
an actual opposing endpoint,
so their required witnesses
are not vacuous.
\end{proof}
\end{annexitem}

\begin{annexitem}{annex:p9.stavi.point.chain.complement}{thm:p09-stavi-point-chain-complement}{The complete labelled-point recursion}
\begin{proof}
We induct on \(m\). At \(m=1\), absence of an
occurrence is exactly \(\operatorname{All}_{\neg P_1}\).
For larger \(m\), if either extreme label
has no occurrence, no required chain exists,
so the first alternative is correct.

If the lower boundary of \(P_1\) is \(a\),
any realizing chain for \(P_2,\ldots,P_m\)
can be preceded by a \(P_1\)-point strictly
between \(a\) and its first point, by
coinitiality. Conversely deleting the
first point of a full chain gives that
shorter chain. Their complements therefore
agree under this boundary condition.
If the lower boundary is an actual
interior \(c\), a full chain's second
point is strictly above \(c\). Conversely,
given a shorter chain above \(c\), use
\(c\) itself as the first point if the
infimum is attained, or choose a
\(P_1\)-point between \(c\) and that
chain's first point if it is unattained.
This proves the corresponding actual
boundary alternative with its retained
descriptor. Reversing the order gives
both upper-boundary alternatives and
retains every label \(P_1,\ldots,P_m\)
in the indicated shortened lists.

If none of these cases applies, the
first-label lower boundary \(c\) and
last-label upper boundary \(d\) are
proper gaps, described by
\(G_{\mathbf P}\). All \(P_1\)-points
in \((a,b)\) lie after \(c\), coinitially there;
all \(P_m\)-points in \((a,b)\) lie before \(d\),
cofinally there.

If \(c\geq d\), a \(P_1\)-point
cannot precede a \(P_m\)-point,
so no full chain exists.
The partition
\(\operatorname{Part}(\neg P_1,\neg P_m;\varnothing)(a,b)\)
describes exactly this ordering of
the two gaps. For \(c\geq d\),
split at \(c\): the nonempty
left side has no \(P_1\),
and the nonempty right side
has no \(P_m\).
Conversely such a partition
puts every \(P_m\)-point
in its first cell and every
\(P_1\)-point in its second.
Thus all the former precede
all the latter, which implies
\(d\leq c\). Nonemptiness
and cofinality of the two
occurrence sets exclude any
exception caused by coincident
extreme occurrences.

If \(c<d\), a point \(z\)
lies strictly between them
exactly when \(B_{\mathbf P}(a,z,b)\)
holds. In one direction choose
an occurrence from each
side using coinitiality and
cofinality. In the other,
a \(P_1\)-point before \(z\)
forces \(c<z\), and a
\(P_m\)-point after \(z\)
forces \(z<d\). Such
an actual \(z\) exists
between distinct gaps.

Fix any such \(z\).
If a full chain exists,
its first and last points
can be replaced so that
the first lies below \(z\)
and the last above \(z\),
without changing its internal
points. For \(m\geq3\),
choose the first point
below the smaller of \(z\)
and its second point, using
coinitial \(P_1\)-occurrences;
choose the last above the
larger of \(z\) and its
penultimate point, using
cofinal \(P_m\)-occurrences.
For \(m=2\), choose these
two points directly from
the two clauses of
\(B_{\mathbf P}\).

The adjusted chain places
\(z\) either strictly between
two successive points, or
at an internal point with
index \(2,\ldots,m-1\).
A gap cut gives a prefix
chain on \((a,z)\) and
a suffix on \((z,b)\).
A point cut gives such
shorter chains and the
middle label \(P_i(z)\).
Conversely each of those
conditions reconstructs
a full increasing chain.
Thus at this fixed \(z\)
negation of the full chain
is exactly the conjunction
\(C_{\mathbf P}\), by
De Morgan's laws.
The induction hypothesis
supplies every displayed
smaller-chain complement.
Quantifying the same \(z\)
with its retained
\(B_{\mathbf P}\) condition
therefore gives the last
alternative exactly.

The boundary classification
proves that these cases
cover every assignment.
Each guarded alternative
implies the required
negation, so overlaps
cause no problem.
All primitives, two-cell
partitions and their
finite positive combinations
are simple. Existential
closure compiles the
displayed point quantifiers.
Every recursive call uses
strictly fewer than \(m\)
labels, so the construction
terminates. Empty intervals
are already covered by
the absence-of-occurrence
alternative. This proves
the full recursion.
\end{proof}
\end{annexitem}

\begin{annexitem}{annex:p9.stavi.bounded.sampling}{thm:p09-stavi-bounded-sampling}{The full four-case bounded truncation proof}
\begin{proof}
We first record a property of the finite-seed formulas:
if \(i\notin O\), \(F_i(c)\) holds and \(D_i\)
holds throughout \([u,c]\), then \(F_i(u)\)
holds. At the terminal index this follows
from \(F_k=D_k\). Otherwise \(F_i\) uses
either \(D_i\mathbin{\mathsf U_s}F_{i+1}\)
or \(D_i\mathbin{\mathsf U_*}F_{i+1}\).
In the first case retain its later witness
and extend its initial \(D_i\)-interval
leftward to \(u\). In the second case
extend the first cell of its two-cell
partition leftward to \(u\).
That cell is unflagged. These operations
prove the claimed property without
changing any future operand's truth.

If a partition of \([a_1,b]\), \(b\leq d\),
exists, take \(a_1\) as the point in
its first cell and choose one point
from each later nonempty cell.
They are strictly increasing and
at most \(b\). The completed
finite-prefix theorem makes \(F_i\)
true throughout its \(i\)-th cell,
so these are the required samples.

Conversely assume the samples exist.
For each \(\ell\), we construct
a partition of \([a_1,b_\ell]\),
where \(b_\ell\leq a_\ell\),
whose cells have labels
\(D_i\land F_i\), \(1\leq i\leq\ell\),
and the original flags restricted
to those indices. At \(\ell=1\)
use the singleton \(\{a_1\}\);
\(F_1\) implies \(D_1\).

Suppose the construction reaches
an endpoint \(b_0\leq a_\ell\).
Its last cell contains \(b_0\),
so \(F_\ell(b_0)\) holds.
Write \(a=a_{\ell+1}>b_0\).
There are four cases.

If both indices are flagged,
the recurrence gives an
immediate successor \(c>b_0\)
with \(F_{\ell+1}(c)\).
Since \(a>b_0\), \(c\leq a\).
Retain the old partition
and add \(\{c\}\).
The empty neighbor gap
leaves no uncovered point,
and the new label follows
from \(F_{\ell+1}(c)\).

If only \(\ell\) is flagged,
there is \(c>b_0\) with
\(F_{\ell+1}(c)\) and
\(D_{\ell+1}\) throughout
\((b_0,c]\). Let
\(b=\min(c,a)\).
Retain the old last
singleton and add
\((b_0,b]\). This cell
is nonempty and unflagged.
For each \(u\) in it,
the property established
above applies to
\(F_{\ell+1}(c)\) and
the all-\(D_{\ell+1}\)
interval \([u,c]\).
Thus the entire new
cell has both required
labels.

If only \(\ell+1\) is flagged,
there is \(c>b_0\) with
\(F_{\ell+1}(c)\) and
\(D_\ell\) throughout
\([b_0,c)\). Again let
\(b=\min(c,a)\).
Extend the old unflagged
last cell by \([b_0,b)\),
and add the singleton
\(\{b\}\). Every added
point of the old cell
satisfies \(F_\ell\)
using the same witness
\(c\) and its remaining
\(D_\ell\)-gap.
At \(b\), \(F_{\ell+1}\)
comes from the witness
if \(b=c\), and from
the given sample if
\(b=a<c\). It also
implies \(D_{\ell+1}\).

If neither is flagged,
the recurrence and the
two-cell theorem give
\([b_0,c]=J_1\cup J_2\),
where \(J_1<J_2\) are
nonempty, with labels
\(D_\ell,F_{\ell+1}\).
Every point of \(J_1\)
also satisfies \(F_\ell\):
truncate that cell at
the point and retain
the same \(J_2\).
Every point of \(J_2\)
has \(D_{\ell+1}\).
If \(c<a\), join
\(J_1\) to the old
last cell, take
\(J_2\) as the new
cell, and end at \(c\).
If \(c\geq a\) and
\(a\in J_2\), make
the same join and
truncate \(J_2\) at
\(a\), ending there.
If \(a\in J_1\),
extend the old cell
only by the part
of \(J_1\) below
\(a\), and use
\(\{a\}\) as the
new cell. Its
\(F_{\ell+1}\) label
is the supplied
sample. Equality
\(a=c\) belongs
to the \(J_2\)
case. These
possibilities exhaust
the location of
the actual sample,
including a missing
cut between the
two cells.

In every case all
earlier cells are
unchanged, the new
endpoint is at most
\(a_{\ell+1}\), and
every cell retains
\(D_i\land F_i\).
Induction gives
the required
partition ending
at most \(a_k\leq d\).
At \(k=1\) the
base alone proves
the assertion,
including \(d=a_1\).
\end{proof}
\end{annexitem}

\begin{annexitem}{annex:p9.stavi.uniform.handlers}{thm:p09-stavi-uniform-handlers}{The independent prefix complements and uniform-cell identities}
\begin{proof}
Pad the initial point \(a\) with a true forced
singleton as specified. A partition of \((a,z]\)
for \(W\) is exactly a partition of \([a,z]\)
for this padded tuple. By the bounded sampling
lemma, such a partition ending before \(b\)
exists exactly when \(F_1(a)\) and
an increasing chain in \((a,b)\) with
labels \(F_2,\ldots,F_{N+1}\) exist.
For one direction choose samples from
the padded realizing cells. For the
other use the last sample as the
bound \(d<b\); the lemma yields an
endpoint at most that sample. The
padded tuple has at least two cells,
so that endpoint is strictly above
\(a\), as the target requires.
Negating this equivalence gives
the displayed \(L_W\), using the
completed plain-chain complement.
It is simple. Reversal gives
the exact right-hand handler
\(R_W\), with every label and
flag restored after the dual
construction.

Suppose \(1\notin O\) and \(D_1\)
holds throughout \((a,b)\).
Any partition of that interval
can be truncated on the left
at a chosen point \(z\) of its
first cell, giving the same
tuple on \([z,b)\).
Conversely, extend the first
cell of such a suffix partition
leftward by \((a,z)\).
Every added point has \(D_1\);
the first cell is unflagged,
so the extension is permitted.
These operations preserve all
other cells and flags.
The original partition is
therefore equivalent to the
existential suffix condition
whose complement is \(R_W\).
Empty open intervals make
both positive conditions
false and both complements
true. The argument for
a uniform last cell is
the order dual and gives
\(L_W\).

The construction of the
two handlers uses only
the finite prefix formulas,
the independently proved
sampling lemma, and
plain-chain negation.
It does not call the
partition-complement
algorithm at any rank.
The flag restrictions
are substantive: extending
a forced singleton can
change its cardinality.
\end{proof}
\end{annexitem}

\begin{annexitem}{annex:p9.stavi.partition.cut}{thm:p09-stavi-partition-cut}{Every point-cut alternative and its converse}
\begin{proof}
Fix \(a<z<b\). In any realizing partition,
there is a unique cell \(i\) containing \(z\).
If it is flagged, that cell is exactly
\(\{z\}\). All earlier cells partition
\((a,z)\) and all later cells partition
\((z,b)\). When the corresponding list
is empty, the corresponding interval is
empty, which is precisely the specified
\(E\) convention.

For an unflagged cell, its portion below
\(z\) is either empty or is a nonempty
last cell of the left partition. This
gives respectively the lists \(1:i-1\)
and \(1:i\). Its portion above \(z\)
likewise gives \(i+1:N\) or \(i:N\).
The point itself must satisfy \(D_i\).
Every inherited flag belongs to an
unchanged cell; the possibly retained
portions of cell \(i\) are unflagged.
This proves the forward implication.

Conversely, choose the alternatives
in \(A_i\). In the flagged case insert
the singleton \(\{z\}\) between the two
partitions. In the unflagged case
join the selected terminal portion
on the left, \(\{z\}\), and the
selected initial portion on the
right. Their union is nonempty,
convex and \(D_i\)-labelled. The
remaining cells preserve their
order, labels and flags. This
proves the reverse implication
for all four empty-or-nonempty
portion choices, including adjacent
endpoints.

For an interior index, both side
lists have length at most \(N-1\).
Explicitly, if \(i\in O\), its
complement is
\[
 N_{1:i-1}(a,z)\lor\neg D_i(z)
                       \lor N_{i+1:N}(z,b).
\]
If \(i\notin O\), it is
\[
 \bigl(N_{1:i-1}(a,z)\land N_{1:i}(a,z)\bigr)
 \lor\neg D_i(z)\lor
 \bigl(N_{i+1:N}(z,b)\land N_{i:N}(z,b)\bigr).
\]
An empty-list complement means
\(\operatorname{Occ}_{\top}\).
These are actual finite positive
combinations of smaller outputs.
Negating the equivalence at the
fixed \(z\), then conjoining any
descriptor \(H(a,z,b)\), is valid
pointwise. Existentially quantifying
that same \(z\) preserves it;
no witness is substituted from a
different boundary condition.
\end{proof}
\end{annexitem}

\begin{annexitem}{annex:p9.stavi.equal.gap.branch}{thm:p09-stavi-equal-gap-branch}{Recognition and partitioning at one common gap}
\begin{proof}
First assume \(G_W(a,b)\). Its first
conjuncts give a unique proper gap
\(c\in(a,b)\), with \(D_1\) throughout
\((a,c)\) and failures of \(D_1\)
coinitial on the right of \(c\).
Its last conjuncts give a unique
proper gap \(d\in(a,b)\), with \(D_N\)
throughout \((d,b)\) and failures of
\(D_N\) cofinal on the left of \(d\).
The bounded occurrence clauses ensure
that both gaps lie strictly between
the actual endpoints.

We first prove
\[
 \operatorname{Part}(\langle D_1,D_N\rangle,\varnothing)(a,b)
                         \quad\Longleftrightarrow\quad d\leq c.
\]
If \(c=d\), split the interval at
that gap. Both parts are nonempty.
If \(d<c\), two distinct cuts have
an actual point \(z\) strictly
between them: otherwise their lower
sets would be equal. The cells
\((a,z]\) and \((z,b)\) have the
required labels and are nonempty.
Conversely, the initial \(D_1\) cell
of any two-cell partition contains
no point above \(c\), since such
a point would force a failure
coinitial above \(c\) into that
convex initial cell. The final
\(D_N\) cell contains no point
below \(d\), by the dual argument.
If \(c<d\), an actual point between
them belongs to neither cell.
This contradiction proves the
equivalence.

Under these descriptors and
the two-cell condition,
\(T_W(a)\) holds exactly when
\(d<c\). Indeed, choose
\(z\in(d,c)\) when \(d<c\).
Then \(D_1\) holds on \((a,z]\);
its gap from \(z\) is the same
\(c\), and \(D_N\) holds on
\((z,c)\). The proved meaning of
\(\operatorname{UG}\) therefore
makes \(z\) an until witness.
For the converse, any until
witness \(z\) must lie below
\(c\): otherwise the \(D_1\)
initial segment encounters a
failure above \(c\). Its
\(D_1\)-gap is consequently
still \(c\). The two-cell
condition already gives
\(d\leq c\). In the equality
case, every point below \(c\)
has a \(D_N\) failure between
it and \(c\), by the cofinal
failure condition at \(d\).
Thus equality
makes \(\operatorname{UG}(D_1,D_N)\)
false at every possible witness.
Together with the two-cell
equivalence this proves exactly
\(Q_W\leftrightarrow(G_W\land c=d)\).
No stronger claim about \(T_W\)
outside the two-cell condition
is required.

Write \(g=c=d\). The open-gap
translation theorem now gives
\[
 \lambda_i(a)\ \longleftrightarrow\
             \operatorname{Part}(\langle D_1,\ldots,D_i\rangle,O_{1:i})(a,g).
\]
Reverse the entire order and
every label in the same theorem.
The suffix ending at \(b\)
becomes a prefix beginning at
\(b\); an original index \(\ell\)
becomes \(N+1-\ell\). Reversing
back gives the displayed
meaning of \(\rho_j(b)\) on
\((g,b)\). The retained \(Q_W\)
ensures that these two independent
descriptions name the same cut.
All temporal evaluations are
at \(a,b\), which are actual
points.

In a finite partition of \((a,b)\),
either some cell has points on
both sides of \(g\), or none
does. In the second case some
unique \(i\), \(1\leq i<N\),
is the last cell on the left.
Both sides of a proper interior
gap are nonempty, so neither
list is empty. Their partitions
give \(\lambda_i\land\rho_{i+1}\).
In the first case there is a
unique crossing cell \(i\).
Its two nonempty portions give
\(\lambda_i\land\rho_i\).
It cannot be flagged, since
a singleton lies on only one
side of a proper gap.

Conversely, a between-cell pair
can be concatenated across \(g\).
For a within-cell pair with
\(i\notin O\), merge its two
portions of cell \(i\).
Their union is convex in the
original order: there is no
point at \(g\), and each
portion reaches the cut.
Every point of the union has
\(D_i\); its size is allowed
by the absence of the flag.
All other cells and flags
are unchanged. These operations
establish that the finite
disjunction of the stated
pairs is equivalent to the
original partition. Its exact
De Morgan complement is
\(Z_W\). Negated \(\lambda_i\)
and \(\rho_j\) are Boolean
negations of already constructed
temporal formulas, hence valid
canonical unary labels.

For the rational example, the
listed three intervals form a
partition and the middle interval
crosses \(g\), so its two
portions witness the within-cell
index \(2\). If this cell is
forced to be a point \(z<g\),
choose a rational strictly
between \(z\) and \(g\);
it lies in the final cell
and fails \(D_3\). If \(z>g\),
a rational between \(g\) and
\(z\) lies in the initial
cell and fails \(D_1\).
There is no third position
for \(z\). Moreover \(D_3\)
fails throughout \((z,g)\)
for every \(z<g\), whereas
\(D_2=\top\) holds there.
This proves both the outcomes
and the asserted necessity
of the final-label index.
\end{proof}
\end{annexitem}

\begin{annexitem}{annex:p9.stavi.open.partition.complement}{thm:p09-stavi-open-partition-complement}{The complete terminating partition-negation induction}
\begin{proof}
All assertions below carry the
guard \(a<b\). For \(N=1\),
an unflagged partition exists
exactly when \((a,b)\) is
nonempty and every point has
\(D_1\). For a flagged cell
it exists exactly when that
interval has one point and
the point has \(D_1\).
Their displayed complements
are therefore exact.

Suppose all shorter words have
their constructed complements.
If the initial cell is forced
to be a singleton, the first
point of \((a,b)\) must be the
immediate successor of \(a\).
If \(a\) has no successor,
no such partition is possible.
If its successor exists, it
is at most \(b\); equality
is exactly the empty-interval
case. Otherwise it is the
unique \(z\in(a,b)\) with
\(E(a,z)\). The original
partition exists exactly when
\(D_1(z)\) and the shorter
word partitions \((z,b)\).
This proves every alternative
of the first forced-cell
formula. The last-cell formula
is the exact order dual.
When both outer cells are
flagged the first procedure
applies, and its shorter
recursive call handles the
remaining flag. Thus forced
cells are removed before
any extension of an outer
cell is attempted.

Assume now \(1,N\notin O\).
Emptiness is immediate.
If \(D_1\) holds throughout
the nonempty interval, the
independent suffix handler
gives exactly its complement.
If \(D_N\) holds throughout,
the prefix handler does so.
Each condition is sound on
its entire stated domain,
including overlapping domains.
If failures of \(D_1\) occur
arbitrarily close to \(a\)
from the right, no initial
nonempty \(D_1\) cell is
possible. To see this, choose
a point of that cell; any
failure strictly between
\(a\) and that point must
belong to the same initial
cell, a contradiction.
The last-cell argument at
\(b\) is dual. These justify
the endpoint-limit branches.

For actual boundary points
we compute the two extreme
point-cut alternatives explicitly.
Since the outer cells are
unflagged, their formulas
simplify to
\[
\begin{split}
 A_1={}&\operatorname{All}_{D_1}(a,z)\land D_1(z)
              \land(\Pi_{2:N}(z,b)\lor\Pi_W(z,b)),\\
 A_N={}&(\Pi_{1:N-1}(a,z)\lor\Pi_W(a,z))
              \land D_N(z)\land\operatorname{All}_{D_N}(z,b).
\end{split}
\]
Here \(E\lor\Pi_{1:1}\) is
\(\operatorname{All}_{D_1}\),
and the last-cell identity
is the same. This also
accounts for adjacent
subinterval endpoints.

Under \(I_W(a,z,b)\), the
left interval is uniformly
\(D_1\). If \(D_1(z)\)
holds, the other disjunct
of the infimum descriptor
is \(\partial^+\neg D_1(z)\).
It prohibits a full-word
partition of \((z,b)\)
by the initial-cell argument.
Consequently
\[
       \neg A_1\ \longleftrightarrow\
                     \neg D_1(z)\lor N_{2:N}(z,b).
\]
For the other extreme, direct
De Morgan expansion gives
\[
 \neg A_N\ \longleftrightarrow\
 \neg D_N(z)\lor\operatorname{Occ}_{\neg D_N}(z,b)
 \lor\bigl(N_{1:N-1}(a,z)\land\neg\Pi_W(a,z)\bigr).
\]
The last full-word complement
is \(R_W(a,z)\), because
\(D_1\) is uniform there
and cell \(1\) is unflagged.
This is exactly the written
\(I_W\) branch, with \(M_W\)
providing all interior indices.
The point-cut equivalence
therefore proves it at the
same \(z\), without any
same-rank inductive premise.

Under \(J_W\), the right
interval is uniformly \(D_N\).
If \(D_N(z)\) holds,
\(\partial^-\neg D_N(z)\)
prohibits a full-word
partition of \((a,z)\).
Thus \(\neg A_N\) becomes
\(\neg D_N(z)\lor N_{1:N-1}(a,z)\).
The other extreme expands as
\[
 \neg A_1\ \longleftrightarrow\
 \neg D_1(z)\lor\operatorname{Occ}_{\neg D_1}(a,z)
 \lor\bigl(N_{2:N}(z,b)\land L_W(z,b)\bigr).
\]
This proves the written
\(J_W\) branch and all its
endpoint-label alternatives.

It remains to consider two
proper failure gaps \(c,d\).
Under \(G_W\), the condition
\(\operatorname{All}_{D_1}(a,z)\land D_1(z)\)
means exactly \(z<c\): below
\(c\) every such point has
the label, whereas above
\(c\) a failure occurs
between \(c\) and \(z\).
Its final-label dual means
exactly \(d<z\). Accordingly
\(V_W\) means \(d<z<c\)
and has an actual witness
exactly when \(d<c\).
Under this condition the
extremes reduce to
\[
\begin{aligned}
 \neg A_1&\longleftrightarrow
                   N_{2:N}(z,b)\land L_W(z,b),\\
 \neg A_N&\longleftrightarrow
                   N_{1:N-1}(a,z)\land R_W(a,z).
\end{aligned}
\]
The replacements use uniform
final and initial cells
respectively. Together with
\(M_W\) these prove the
overlap-gap branch.

Similarly,
\(\operatorname{Occ}_{\neg D_1}(a,z)\)
means \(c<z\), and
\(\operatorname{Occ}_{\neg D_N}(z,b)\)
means \(z<d\). Thus \(B_W\)
means \(c<z<d\), with
a witness exactly when
\(c<d\). At such a point
\(A_1\) and \(A_N\) are
already false: the first
would require a uniform
left \(D_1\) interval,
the last a uniform right
\(D_N\) interval. Hence
the exact complement is
\(M_W\), giving the
between-gap branch.
When \(c=d\), the preceding
common-gap lemma proves
the remaining \(Q_W\land Z_W\)
alternative.

For exhaustion, if the
interval is neither empty
nor uniform for either
outer label, both failure
sets are nonempty. Classify
the lower boundary of the
first set and the upper
boundary of the last set
by the elementary boundary
lemma. An outer endpoint
is handled by its limit
branch. An actual interior
point is handled by
\(I_W\) or \(J_W\).
Otherwise both boundaries
are proper gaps and exactly
one of \(d<c,c<d,c=d\)
holds. The three gap
constructions apply. Thus
every failure belongs
to at least one branch,
and every branch has
already been proved to
imply failure.

Finally, construct all
plain-chain complements first,
then the bounded-prefix and
bounded-suffix handlers, and
then \(N_W\) by word length.
Every occurrence of \(N\)
on the right has fewer
cells. The same-length
\(L_W,R_W\) calls use
the separate chain recursion;
\(\lambda_i,\rho_j\) use
the proved finite gap
translator. None calls
the present partition
induction. The elementary
conditions, order atoms,
modal labels, finite
conjunctions and disjunctions,
and existential projections
are compiled by positive
partition closure. Empty
conjunctions are true and
empty disjunctions false.
The resulting simple formula
is finite and has no
untranslated negative
partition atom. Conjoining
\(a<b\) gives the asserted
guarded equivalence; adding
\(b<a\lor a=b\) gives its
unguarded complement.
\end{proof}
\end{annexitem}

\begin{annexitem}{annex:p9.stavi.simple.negation}{thm:p09-stavi-simple-negation}{Negating every basic formula and its finite assembly}
\begin{proof}
For order atoms use the strict-order
trichotomy:
\[
 \neg(a=b)\equiv(a<b)\lor(b<a),\qquad
 \neg(a<b)\equiv(a=b)\lor(b<a).
\]
Unary temporal labels are closed
under Boolean negation. The
remaining basic atoms are
the partition formulas on
closed bounded intervals,
rays and the whole order.

Consider \(W\) on \([a,b]\).
If \(b<a\), the interval
is empty, so a positive-length
partition is false. If \(a=b\),
it is true exactly when
\(N=1\) and \(D_1(a)\);
a singleton satisfies either
choice of its flag.
These cases have the displayed
order and unary-label
presentations.

Suppose \(a<b\). If \(N=1\)
and the cell is flagged,
the closed interval contains
the two distinct endpoints
and cannot be a singleton.
If it is unflagged, its
partition condition is
\[
 D_1(a)\land D_1(b)\land
                           \operatorname{All}_{D_1}(a,b).
\]
Its complement is
\(\neg D_1(a)\lor\neg D_1(b)
       \lor\operatorname{Occ}_{\neg D_1}(a,b)\).
The open interval may be
empty here; the endpoint
labels still apply.

For \(N\geq2\), every closed
partition has \(a\) in its
first cell and \(b\) in
its distinct last cell.
It therefore requires
\(D_1(a)\land D_N(b)\).
Remove those endpoints.
For each
\((e_L,e_R)\in\{0,1\}^2\),
interpret \(e_L=1\) as
a nonempty remainder of
the first cell and
\(e_R=1\) as a nonempty
remainder of the last.
Forbid \(e_L=1\) if
\(1\in O\), and forbid
\(e_R=1\) if \(N\in O\).
The reduced word \(W_e\)
retains all middle cells
\(2,\ldots,N-1\), their
flags shifted to their
new positions, and retains
the first or last label
exactly when its \(e\)
is \(1\). Any retained
endpoint cell is unflagged.
It is permissible for the
reduced list to be empty;
this occurs for \(N=2\)
and \(e_L=e_R=0\).
In that case its meaning
is \(E(a,b)\).

Restriction of a realizing
closed partition gives
exactly one of these
open reduced partitions.
Conversely, add \(a\)
to the first retained
cell if \(e_L=1\), or
insert \(\{a\}\) if
\(e_L=0\); perform the
dual operation at \(b\).
The endpoint labels
ensure the enlarged cells
are correctly labelled.
Their convexity follows
from being the first
and last cells of the
whole open interval.
The forbidden flag choices
ensure every forced cell
remains a singleton.
When the reduced list
is empty, \(E\) supplies
exactly the interval
consisting of the two
endpoints. Hence
\[
 \operatorname{Part}(W)[a,b]\ \longleftrightarrow\
 D_1(a)\land D_N(b)\land
                  \bigvee_{e\text{ allowed}}\Pi_{W_e}(a,b).
\]
Its complement is the
disjunction of the two
endpoint-label failures
and
\[
                 \bigwedge_{e\text{ allowed}}N_{W_e}(a,b).
\]
For an empty reduced
word use
\(\operatorname{Occ}_{\top}(a,b)\).
All open-word complements,
including length \(N\)
when both remainders are
nonempty, have already
been constructed. Thus
this step is not a
same-rank recursive call
in the open-word proof.

A ray partition with its
one actual free endpoint
has an explicit formula
\(A\in\mathcal T_{\mathrm{gap}}\)
by the basic temporal
translation theorem.
Its complement is the
unary label \(\neg A\),
represented as the
one-cell partition
\(\operatorname{Part}(\langle\neg A\rangle,\varnothing)[x,x]\).
For a whole-order
partition, that theorem
likewise supplies a
temporal formula \(A\)
whose value at every
actual point is exactly
the same global partition
sentence. Thus its
complement is represented
by the whole-order
one-cell partition
\(\operatorname{Part}(\langle\neg A\rangle,\varnothing)\).
The order is nonempty,
and \(A\)'s global
meaning makes \(\neg A\)
uniformly true precisely
when that original
sentence fails. This
also treats singleton
orders and rays with
an actual first or last
point; the existing ray
translation already
retains those cases.
Formal infinite endpoints
appear only as interval
notation. They are never
arguments of \(A\).

Finally a simple formula
is a finite disjunction
of basic conjunctions.
Apply the two De Morgan
laws, replace each
negated basic atom by
the construction above,
and distribute the
resulting finite positive
combinations through the
proved intersection and
union construction.
An empty conjunction
negates to false, and
an empty disjunction
negates to true. Both
constants already have
simple presentations.
This gives effective
negation of every simple
formula with any finite
list of free variables.
\end{proof}
\end{annexitem}

\begin{annexitem}{annex:p9.stavi.completeness}{thm:p09-stavi-completeness}{The full first-order induction and temporal return}
\begin{proof}
Work with a fixed finite
monadic first-order input
and retain its finite list
of free variables at every
stage. An atomic unary
formula \(P(x)\) is the
singleton-interval partition
\(\operatorname{Part}(\langle P\rangle,\varnothing)[x,x]\).
Equality and strict-order
atoms are basic partition
atoms. Truth and falsity
use their simple constants.
This starts the induction.

Given the constructed
simple formulas for two
inputs, the positive
closure theorem constructs
their conjunction by its
finite common refinements
and their disjunction by
union of the finite lists.
It retains every weak
order type of shared
free variables, including
equalities, and every
singleton flag in the
refinement. For negation
apply the completed
simple-negation theorem.
For \(\exists x\,\varphi\),
apply the positive
existential-elimination
construction to the
already obtained simple
formula for \(\varphi\).
That construction first
anchors the named actual
points as singleton
cells in their weak
order types. If \(x\)
equals a retained point,
conjoin its labels at
that point and remove
only its name. If it
is a distinct point,
erase its name from
the anchored partition,
retaining the singleton
cell and its label;
existence of that cell
is exactly the former
existential witness.
The adjacent cell lists
are glued in their
finite refinement order.
Surviving strict bounds
remain actual distinct
endpoints. The two
unbounded outer regions
remain rays, and no
witness is placed at
a formal infinite
endpoint. These are
precisely the cases
proved in
\cref{annex:p9.stavi.positive.partitions}.
Universal quantification
is implemented as
\(\neg\exists x\neg\);
the two required
negation operations
are now available.
This finishes the
structural induction
for arbitrary finite
free-variable lists.

For an output with
only \(x\) free,
apply the full
one-variable simple
partition translation.
Each closed basic
interval has its
endpoints identified
with \(x\), each
one-sided interval
has the proved
finite/ray prefix
translation or its
past dual, and the
whole-order condition
has its anywhere
translation. Finite
Boolean combinations
are preserved. The
basic translation's
truth at \(x\) is
evaluated in the
original whole model:
no temporal label
is implicitly
relativized to the
cell on which it
is used. In a
sentence case the
translation has the
same truth value
at every actual
point, and the
nonempty model
provides such a
point.

Here is the actual final
return on each basic form.
The atom \(x=x\) returns
\(\top\), and \(x<x\)
returns \(\bot\).
The closed partition
\(\operatorname{Part}(W)[x,x]\)
returns \(D_1\) if \(N=1\),
and \(\bot\) if \(N>1\).
A right ray returns the
ray-seed formula \(F_1\)
of the four-case recurrence;
a left ray returns that
same construction with
the tuple and flags
reversed and each modality
dualized. For the whole
order, let \(F_W\) denote
that right-ray formula,
and define the actual
one-cell past formula
\[
 H_1=
 D_1\land
 \begin{cases}
  \neg(\top\mathbin{\mathsf S_s}\top),&1\in O,\\
  \neg(\top\mathbin{\mathsf S_s}\neg D_1),&1\notin O.
 \end{cases}
\]
Return
\[
 A\lor(\top\mathbin{\mathsf U_s}A)
       \lor(\top\mathbin{\mathsf S_s}A),
             \qquad A=H_1\land F_W.
\]
At a witness for \(A\),
the one-cell past interval
and the first future cell
meet at that actual point.
Their union is the
original first cell;
the past flag condition
ensures it remains a
singleton when required.
The other future cells
are unchanged. Conversely,
any point of the first
cell of a whole-order
partition gives such
a witness. Strict
future, strict past
and the present
disjunct exhaust every
possible witness relative
to the evaluation point.
This proves the formula
also at first or last
points and on a
singleton order.
Finite conjunction and
disjunction of these
returns gives the
specified simple formula's
temporal return.

The canonical expansion
is effective and
introduces no new
expressive premise.
Each new unary name
used above abbreviates
a displayed finite
formula of
\(\mathcal T_{\mathrm{gap}}\)
over earlier labels:
Boolean combinations,
the four primitive
modalities, the
derived limit and
gap formulas, the
finite prefix
recurrences, or the
finite gap translators.
Plain-chain negation
decreases its number
of labels. The
open-word induction
decreases its number
of cells except for
its independently
constructed handlers.
Every refinement and
distribution enumerates
a finite set determined
by its finite input.
Finally the outer
first-order induction
decreases syntax-tree
size. Consequently
all intermediate
formulas and names
are produced by
terminating explicit
finite operations.
Inline the names
in their creation
order. This returns
an actual finite
formula in the
four stated primitive
modalities over the
original unary
predicate alphabet.

For the converse
translate atoms and
Boolean operations
homomorphically.
At a strict-until
constructor use the
proved existential
future witness and
universal open
intermediate-interval
clause; at a strict-since
constructor use its
order reverse.
At a gap-until
constructor use the
first-order formula
of
\cref{thm:p09-stavi-first-order-clauses},
including its
definable-gap condition
and its bounded
right-hand label
condition. At gap-since
use the exact
reversed clause.
Replace each predicate
placeholder by the
recursively obtained
operand translation,
with capture-free
fresh bound variables.
The truth-table
equivalences and
syntax induction
prove preservation
at every actual
point. No gap
is quantified as
an extra domain
element.

All constructions used
only strict linearity,
the actual existence
conditions explicitly
shown, and finiteness
of the formulas.
In particular they
did not assume
Dedekind completeness,
density, local
discreteness, or
endpoints. The
two translations
therefore establish
the stated equivalence
on every nonempty
linear order and
every interpretation
of its unary
predicates.
\end{proof}
\end{annexitem}
